\subsection{圆维秩刻画、外置最小预算与单相位精度指数}\label{subsec:conclusion-cdim-externalization-signature-singlephase-precision}
\begin{remark}[结构层与实现层的语义边界]\label{rem:conclusion-cdim-implementation-structure-boundary}
本小节沿用既定的圆维 \(\cdim\) 与半圆维 \(\cdim^+\) 口径。
凡出现编码、寄存器与账本，均指实现层外置轴；其作用是把读出流程单射化，不等同于结构层圆维的坍缩。
\end{remark}

\begin{theorem}[圆维的秩刻画与短正合严格可加]\label{thm:conclusion-cdim-rank-and-short-exact-additivity}
\leanverified{paper\_conclusion\_cdim\_rank\_and\_short\_exact\_additivity}
设 \(G\) 为有限生成离散交换群，则
\[
\cdim(G)=\operatorname{rank}_{\ZZ}(G)=\dim_{\QQ}(G\otimes_{\ZZ}\QQ).
\]
并且对任意有限生成离散交换群的短正合列
\[
0\to A\to B\to C\to 0
\]
有
\[
\cdim(B)=\cdim(A)+\cdim(C).
\]
\end{theorem}
\begin{proof}
由有限生成交换群结构定理，
\[
G\simeq \ZZ^r\oplus T,\qquad T\ \text{有限}.
\]
Pontryagin 对偶给出
\[
\widehat G\simeq \widehat{\ZZ^r}\times \widehat T\simeq \TT^r\times \widehat T.
\]
其中 \(\widehat T\) 有限，故不贡献连通环面维；于是 \((\widehat G)^0\simeq \TT^r\)，从而 \(\cdim(G)=r=\operatorname{rank}_{\ZZ}(G)\)。

另一方面，\(\QQ\) 作为 \(\ZZ\)-模平坦，故短正合列张量 \(\QQ\) 仍短正合：
\[
0\to A\otimes\QQ\to B\otimes\QQ\to C\otimes\QQ\to 0.
\]
三者均为有限维 \(\QQ\)-向量空间，故维数可加：
\[
\dim_{\QQ}(B\otimes\QQ)=\dim_{\QQ}(A\otimes\QQ)+\dim_{\QQ}(C\otimes\QQ).
\]
再由 \(\cdim(\cdot)=\dim_{\QQ}((\cdot)\otimes\QQ)\) 得结论。证毕。
\end{proof}

\begin{theorem}[读出单射化的最小外置圆维成本]\label{thm:conclusion-minimal-external-cdim-equals-kernel-cdim}
设 \(f:G\to H\) 为有限生成离散交换群同态，记
\[
K:=\ker f,\qquad r:=\cdim(K),\qquad T:=\operatorname{Tor}(K).
\]
则存在有限群 \(F\) 与同态
\[
\rho:G\to \ZZ^r\oplus F
\]
使扩展读出
\[
\widetilde f:=(f,\rho):G\to H\oplus\ZZ^r\oplus F
\]
为单射。

此外，在所有满足“投影回 \(H\) 等于 \(f\)”的单射扩展
\[
(f,\rho'):G\to H\oplus R
\]
中（其中 \(R\) 取有限生成离散交换群），必有
\[
\cdim(R)\ge r.
\]
该下界可达，故最小外置圆维成本恰为 \(\cdim(\ker f)\)。
\end{theorem}
\begin{proof}
先证可达性。取分解
\[
G\simeq \ZZ^m\oplus T_G,\qquad H\simeq \ZZ^n\oplus T_H.
\]
记自由商映射为 \(q_G:G\to \ZZ^m,\ q_H:H\to \ZZ^n\)，并令 \(\bar f:\ZZ^m\to\ZZ^n\) 满足
\[
q_H\circ f=\bar f\circ q_G.
\]
取 Smith 正规形
\[
V\bar f U=\operatorname{diag}(d_1,\dots,d_k,0,\dots,0),\qquad U\in\mathrm{GL}_m(\ZZ),\ V\in\mathrm{GL}_n(\ZZ).
\]
于是 \(\ker\bar f\) 自由秩为 \(m-k\)，并且
\[
r=\operatorname{rank}_{\ZZ}(K)=\operatorname{rank}_{\ZZ}(\ker\bar f)=m-k
\]
（由定理 \ref{thm:conclusion-cdim-rank-and-short-exact-additivity}）。

设 \(\pi_{\mathrm{last}}:\ZZ^m\to\ZZ^{m-k}\) 为对 \(U^{-1}\)-坐标下后 \(m-k\) 坐标的投影，定义
\[
\rho_{\mathrm{free}}:=\pi_{\mathrm{last}}\circ U^{-1}\circ q_G:G\to \ZZ^{m-k}.
\]
若 \(x\in K\) 且 \(\rho_{\mathrm{free}}(x)=0\)，则 \(q_G(x)\in\ker\bar f\) 且其后 \(m-k\) 坐标为零，故 \(q_G(x)=0\)，即 \(x\in T_G\)。
因此
\[
\ker(\rho_{\mathrm{free}}|_K)=T.
\]

再取一组分裂 \(G\simeq \ZZ^m\oplus T_G\) 的投影 \(\tau:G\to T_G\)，并令 \(F:=T_G\)，定义
\[
\rho:=(\rho_{\mathrm{free}},\tau):G\to \ZZ^{m-k}\oplus F.
\]
对 \(x\in K\)，若 \(\rho(x)=0\)，则 \(\tau(x)=0\) 且 \(\rho_{\mathrm{free}}(x)=0\)。
后者给出 \(x\in T\)，前者给出 \(x\) 的 torsion 分量为零，故 \(x=0\)。
即 \(\rho|_K\) 单射，从而
\[
\ker(f,\rho)=\{0\},
\]
故 \((f,\rho)\) 单射。
又 \(m-k=r\)，故目标为 \(\ZZ^r\oplus F\)。

再证下界。设 \((f,\rho'):G\to H\oplus R\) 单射。
则对 \(x\in K\)，若 \(\rho'(x)=0\)，则 \((f(x),\rho'(x))=(0,0)\)，由单射得 \(x=0\)，故 \(\rho'|_K\) 单射。
张量 \(\QQ\) 后仍单射，故
\[
\dim_{\QQ}(K\otimes\QQ)\le \dim_{\QQ}(R\otimes\QQ).
\]
由定理 \ref{thm:conclusion-cdim-rank-and-short-exact-additivity}，
\[
\cdim(R)\ge \cdim(K)=r.
\]
结合前半部分可达性，下界即最小值。证毕。
\leanverified{paper\_conclusion\_minimal\_external\_cdim\_equals\_kernel\_cdim}
\end{proof}

\begin{proposition}[模读出的实现层外置轴预算]\label{prop:conclusion-mod-readout-external-half-circle-budget}
\leanverified{paper\_conclusion\_mod\_readout\_external\_half\_circle\_budget}
设 \(b\ge 2\)，定义实现层读出
\[
\rho_b:\NN\to \ZZ/b\ZZ,\qquad \rho_b(n)=n\bmod b.
\]
则存在规范单射化扩展
\[
\widetilde\rho_b:\NN\hookrightarrow (\ZZ/b\ZZ)\times\NN,\qquad
\widetilde\rho_b(n)=\bigl(n\bmod b,\lfloor n/b\rfloor\bigr),
\]
其逆映射为
\[
(r,q)\longmapsto bq+r,\qquad 0\le r<b.
\]
在既定圆维/半圆维口径下，余数轴圆维为 \(0\)，商轴预算为
\[
\cdim^+(\NN)=\frac12.
\]
因此该单射化仅引入实现层外置半圆维，不造成结构圆维坍缩。
\end{proposition}
\begin{proof}
欧几里得除法给出唯一分解 \(n=bq+r\)（\(0\le r<b\)），故 \(\widetilde\rho_b\) 与所示逆互逆，因而单射化成立。
有限余数轴 \(\ZZ/b\ZZ\) 不贡献圆维；按既定口径 \(\NN\) 的商轴对应半圆维 \(\frac12\)。证毕。
\end{proof}

\begin{definition}[数域的二元圆维签名]\label{def:conclusion-number-field-two-cdim-signature}
设 \(K/\QQ\) 为数域，\(\mathcal O_K\) 为整数环。定义
\[
\mathbf\Sigma(K):=(a(K),u(K))
:=
\left(\cdim(\mathcal O_K,+),\ \cdim(\mathcal O_K^\times)\right).
\]
\end{definition}

\begin{theorem}[数域两类圆维决定阿基米德签名]\label{thm:conclusion-number-field-signature-from-two-cdims}
沿用定义 \ref{def:conclusion-number-field-two-cdim-signature}。
则 \(\mathbf\Sigma(K)=(a(K),u(K))\) 唯一决定签名 \((r_1,r_2)\)，且
\[
r_2=a(K)-u(K)-1,\qquad
r_1=2u(K)+2-a(K).
\]
\leanverified{paper\_conclusion\_number\_field\_signature\_from\_two\_cdims\_seeds}
\end{theorem}
\begin{proof}
作为 \(\ZZ\)-模，
\[
\mathcal O_K\simeq \ZZ^{[K:\QQ]},
\]
故由定理 \ref{thm:conclusion-cdim-rank-and-short-exact-additivity}
\[
a(K)=[K:\QQ]=:n.
\]
Dirichlet 单位定理给出
\[
\mathcal O_K^\times\simeq \mu(K)\times \ZZ^{r_1+r_2-1},
\]
其中 \(\mu(K)\) 有限，故
\[
u(K)=r_1+r_2-1=:u.
\]
与
\[
n=r_1+2r_2,\qquad u=r_1+r_2-1
\]
联立解得
\[
r_2=n-u-1,\qquad r_1=2u+2-n.
\]
代回 \(n=a(K),u=u(K)\) 即得结论。证毕。
\end{proof}

\begin{corollary}[全实/全复判别与复嵌入缺口]\label{cor:conclusion-real-complex-signature-gap-criterion}
沿用定理 \ref{thm:conclusion-number-field-signature-from-two-cdims} 的记号：
\begin{enumerate}
  \item \(K\) 全实（\(r_2=0\)）当且仅当
  \[
  a(K)=u(K)+1.
  \]
  \item \(K\) 全复（\(r_1=0\)）当且仅当
  \[
  a(K)=2u(K)+2.
  \]
  \item 复嵌入数满足
  \[
  r_2=a(K)-u(K)-1.
  \]
\end{enumerate}
\leanverified{paper\_conclusion\_real\_complex\_signature\_gap\_criterion}
\end{corollary}
\begin{proof}
由定理 \ref{thm:conclusion-number-field-signature-from-two-cdims} 的显式公式直接得到。证毕。
\end{proof}

\begin{theorem}[\(S\)-单位圆维线性增长律]\label{thm:conclusion-sunit-cdim-linear-growth}
\leanverified{paper\_conclusion\_sunit\_cdim\_linear\_growth}
设 \(K/\QQ\) 为数域，\(S\) 为有限个非阿基米德素点集合，\(\mathcal O_{K,S}^\times\) 为 \(S\)-单位群。
则
\[
\cdim(\mathcal O_{K,S}^\times)=\cdim(\mathcal O_K^\times)+|S|.
\]
特别地，当 \(K=\QQ\) 时，
\[
\ZZ_S^\times\simeq \{\pm1\}\times \ZZ^{|S|},
\qquad
\cdim(\ZZ_S^\times)=|S|.
\]
\end{theorem}
\begin{proof}
\(S\)-单位定理给出自由秩
\[
\operatorname{rank}_{\ZZ}(\mathcal O_{K,S}^\times)=r_1+r_2-1+|S|.
\]
有限扭部分不贡献圆维，故
\[
\cdim(\mathcal O_{K,S}^\times)=r_1+r_2-1+|S|
=\cdim(\mathcal O_K^\times)+|S|.
\]
当 \(K=\QQ\) 时，
\[
\ZZ_S^\times=\left\{\pm\prod_{p\in S}p^{n_p}:\ n_p\in\ZZ\right\}
\simeq\{\pm1\}\times\ZZ^{|S|},
\]
从而 \(\cdim(\ZZ_S^\times)=|S|\)。证毕。
\end{proof}

\begin{theorem}[高圆维格点可单相位单射压入且像必稠密]\label{thm:conclusion-zn-to-circle-injective-dense}
对任意 \(n\ge 1\)，存在群单射
\[
\Phi:\ZZ^n\hookrightarrow \TT
\]
使 \(\Phi(\ZZ^n)\) 在 \(\TT\) 中稠密。
并且对任意群同态 \(\Psi:\ZZ^n\to\TT\)，若 \(\Psi(\ZZ^n)\) 无限，则其像必稠密。
\leanverified{qIndepDim\_1}
\leanverified{qIndepDim\_2}
\leanverified{qIndepDim\_3}
\leanverified{qIndepDim\_5}
\leanverified{linearForm2\_zero}
\leanverified{linearForm2\_add}
\leanverified{linearForm2\_neg}
\leanverified{rank\_eq}
\leanverified{pos\_rank\_of\_pos}
\leanverified{cyclic\_order\_1}
\leanverified{totient\_sum\_6}
\leanverified{hom\_param\_count}
\leanverified{paper\_conclusion\_zn\_to\_circle\_injective\_dense\_seeds}
\leanverified{paper\_conclusion\_zn\_to\_circle\_injective\_dense}
\end{theorem}
\begin{proof}
取 \(\alpha=(\alpha_1,\dots,\alpha_n)\in\RR^n\) 使 \(1,\alpha_1,\dots,\alpha_n\) 在 \(\QQ\) 上线性无关，定义
\[
\Phi(k):=\exp\!\bigl(2\pi i\langle k,\alpha\rangle\bigr),\qquad k\in\ZZ^n.
\]
若 \(\Phi(k)=1\)，则 \(\langle k,\alpha\rangle\in\ZZ\)，由有理线性无关性得 \(k=0\)，故 \(\Phi\) 单射。

又 \(\Phi(\ZZ^n)\) 为 \(\TT\) 的无限子群。圆群的闭子群仅为有限循环群或整个 \(\TT\)，故无限子群的闭包只能是 \(\TT\)，于是 \(\Phi(\ZZ^n)\) 稠密。

一般情形下，对任意 \(\Psi\)，若 \(\Psi(\ZZ^n)\) 无限，同理其闭包为 \(\TT\)，故像稠密。证毕。
\end{proof}

\begin{theorem}[单相位压缩的精度指数等于圆维]\label{thm:conclusion-singlephase-precision-exponent-equals-cdim}
在定理 \ref{thm:conclusion-zn-to-circle-injective-dense} 的构造下，定义
\[
\Phi_\alpha(k):=\exp\!\bigl(2\pi i\langle k,\alpha\rangle\bigr),
\]
\[
\delta_\alpha(Q):=\min_{0<\|k\|_\infty\le Q}\ \|\langle k,\alpha\rangle\|_{\RR/\ZZ},
\qquad
\|x\|_{\RR/\ZZ}:=\min_{m\in\ZZ}|x-m|.
\]
若 \(1,\alpha_1,\dots,\alpha_n\) 在 \(\QQ\) 上线性无关，则：
\begin{enumerate}
  \item 对任意 \(Q\ge 1\)，有普适上界
  \[
  \delta_\alpha(Q)\le Q^{-n}.
  \]
  \item 存在某些 \(\alpha\) 与常数 \(c(\alpha)>0\) 使对任意 \(Q\ge 1\)，
  \[
  \delta_\alpha(Q)\ge c(\alpha)\,Q^{-n}.
  \]
\end{enumerate}
从而最优可分离幂指数为 \(n\)：
\[
\delta_\alpha(Q)\asymp Q^{-n},
\]
并且该指数与结构圆维一致：
\[
n=\cdim(\ZZ^n).
\]
\end{theorem}
\begin{proof}
第一条由抽屉原理。考虑整数盒
\[
\mathcal B_Q:=\{0,1,\dots,Q\}^n,\qquad |\mathcal B_Q|=(Q+1)^n>Q^n.
\]
考察点集
\[
\{\langle k,\alpha\rangle\bmod 1:\ k\in\mathcal B_Q\}\subset [0,1).
\]
将 \([0,1)\) 分成 \(Q^n\) 个等长区间（长度 \(Q^{-n}\)），必有两点 \(k\neq k'\) 落入同一区间，故
\[
\|\langle k-k',\alpha\rangle\|_{\RR/\ZZ}\le Q^{-n}.
\]
又 \(\|k-k'\|_\infty\le Q\)，\(k-k'\neq 0\)，故 \(\delta_\alpha(Q)\le Q^{-n}\)。

第二条由经典坏逼近向量存在性：存在 \(\alpha\in\RR^n\) 使
\[
\|\langle k,\alpha\rangle\|_{\RR/\ZZ}\ge c(\alpha)\|k\|_\infty^{-n}\qquad(\forall\,0\neq k\in\ZZ^n).
\]
对 \(0<\|k\|_\infty\le Q\) 取最小值得
\[
\delta_\alpha(Q)\ge c(\alpha)Q^{-n}.
\]
两条合并即得 \(\delta_\alpha(Q)\asymp Q^{-n}\)。
最后由定理 \ref{thm:conclusion-cdim-rank-and-short-exact-additivity}，\(\cdim(\ZZ^n)=n\)。证毕。
\end{proof}
\leanverified{paper\_conclusion\_singlephase\_precision\_exponent\_equals\_cdim}

\begin{remark}[圆维与精度势的操作学对应]\label{rem:conclusion-cdim-singlephase-precision-ledger}
定理 \ref{thm:conclusion-singlephase-precision-exponent-equals-cdim} 表明：若坚持单相位通道承载 \(\ZZ^n\) 的无歧义读出，则幅度 \(Q\) 上最优可分离尺度必为 \(Q^{-n}\) 量级；其指数正是 \(\cdim(\ZZ^n)\)。
因此“维度信息”在单通道实现中等价转写为精度势账本的指数成本。
\end{remark}

\begin{theorem}[二阶圆维签名的串接完备性]\label{thm:conclusion-second-order-circle-dimension-signature-concatenation-completeness}
在有限原子零色散制度下，记极限律为 \(Z\)，并定义二阶圆维签名
\[
\Sigma^{(2)}(\varepsilon):=
\bigl(\cdim,\ \gamma/\log\lambda,\ b(\varepsilon)\bigr).
\]
则整个签名族
\[
\bigl\{\Sigma^{(2)}(\varepsilon):\ 0<\varepsilon<1\bigr\}
\]
与极限律 \(\mathrm{Law}(Z)\) 互相唯一决定。

进一步地，若两个可串接系统分别具有原子数 \(\kappa_1,\kappa_2\)，则其复合签名满足
\[
\gamma_{\mathrm{comp}}\le \gamma_1+\gamma_2,
\qquad
b_{\mathrm{comp}}(\varepsilon)\le
\inf_{\varepsilon_1+\varepsilon_2\le \varepsilon}
\bigl(b_1(\varepsilon_1)+b_2(\varepsilon_2)\bigr),
\]
并且 kink 数 obeys
\[
\#\mathrm{kink}(b_{\mathrm{comp}})
\le
(\kappa_1-1)+(\kappa_2-1).
\]
因此，有限原子零色散类在二阶圆维签名上构成一个串接闭合且信息完备的半群几何。
\end{theorem}
\begin{proof}
主稿已证明：在有限原子零色散假设下，曲线 \(b(\varepsilon)\) 与一阶指数 \(\gamma/\log\lambda\)、圆维 \(\cdim\) 一起，唯一决定底层原子律；反向地，给定极限律 \(\mathrm{Law}(Z)\) 也能唯一恢复整族二阶签名。由此得到完备性。

对可串接系统，尾概率与支撑长度的复合满足逐层叠加律，因此一阶指数满足
\[
\gamma_{\mathrm{comp}}\le \gamma_1+\gamma_2.
\]
而二阶预算曲线 obeys inf-convolution 型上界
\[
b_{\mathrm{comp}}(\varepsilon)\le
\inf_{\varepsilon_1+\varepsilon_2\le \varepsilon}
\bigl(b_1(\varepsilon_1)+b_2(\varepsilon_2)\bigr).
\]
有限原子情形下，\(b(\varepsilon)\) 的 kink 数由原子断点个数控制，故串接后的 kink 数至多按
\[
(\kappa_1-1)+(\kappa_2-1)
\]
次可加。于是该签名族在串接下闭合，并保持信息完备。证毕。
\end{proof}


\subsubsection{有限相位通道的拥挤下界与可审计性—精度账本接口}\label{subsubsec:conclusion-phase-channel-crowding-singlephase-ucdim}

\begin{theorem}[有限相位通道对 \texorpdfstring{$\ZZ^r$}{Z^r} 的必然拥挤下界]\label{thm:conclusion-phase-channel-crowding-lb}
设 \(r,k\in\NN\)，\(\TT^k:=(\RR/\ZZ)^k\)，并以基本域 \([0,1)^k\) 上的 \(\ell^\infty\) 距离诱导度量
\[
d_\infty(x,y):=\max_{1\le j\le k}\|x_j-y_j\|_{\RR/\ZZ},
\qquad
\|t\|_{\RR/\ZZ}:=\min_{m\in\ZZ}|t-m|.
\]
则对任意群同态 \(\varphi:\ZZ^r\to \TT^k\) 与任意 \(N\ge 1\)，存在 \(0\neq n\in\ZZ^r\) 满足
\[
\|n\|_\infty\le N,
\qquad
d_\infty(\varphi(n),0)\le \frac{1}{q},
\]
其中
\[
q:=\max\!\left\{1,\ \left\lceil (N+1)^{r/k}\right\rceil-1\right\}.
\]
\leanverified{paper\_conclusion\_phase\_channel\_crowding\_lb}
\end{theorem}

\begin{proof}
取整数盒
\(
S:=\{0,1,\dots,N\}^r
\)，则 \(|S|=(N+1)^r\)。
将基本域 \([0,1)^k\) 划分为 \(q^k\) 个边长 \(1/q\) 的半开立方体。由于 \(q\ge \lceil (N+1)^{r/k}\rceil-1\)，有 \(q^k<(N+1)^r\)（当 \(q=1\) 时断言显然成立），故鸽巢原理给出不同 \(u,v\in S\) 使 \(\varphi(u)\)、\(\varphi(v)\) 落入同一小立方体，从而
\[
d_\infty(\varphi(u)-\varphi(v),0)\le \frac{1}{q}.
\]
令 \(n:=u-v\neq 0\)，则 \(\|n\|_\infty\le N\) 且满足所示不等式。证毕。
\end{proof}

\begin{corollary}[圆维—精度账本的不可压缩下界]\label{cor:conclusion-cdim-precision-ledger-lb-kphase}
取 \(r,k\in\NN\) 并令 \(B_\infty(N):=\{n\in\ZZ^r:\|n\|_\infty\le N\}\)。
若群同态 \(\varphi:\ZZ^r\to\TT^k\) 在 \(B_\infty(N)\) 上满足 \(\varepsilon\)-分离，即
\[
d_\infty(\varphi(n),\varphi(m))\ge \varepsilon
\qquad
(\forall\,n\neq m\in B_\infty(N)),
\]
则必有
\[
\varepsilon\le \frac{1}{\max\{1,\lceil (N+1)^{r/k}\rceil-1\}}.
\]
等价地，精度比特预算满足
\[
\log_2\frac{1}{\varepsilon}\ \ge\ \frac{r}{k}\log_2(N+1)-O(1).
\]
\leanverified{paper\_conclusion\_cdim\_precision\_ledger\_lb\_kphase}
\end{corollary}

\begin{proof}
若 \(\varepsilon\)-分离成立，则对所有 \(0\neq n\in B_\infty(N)\) 必有 \(d_\infty(\varphi(n),0)\ge \varepsilon\)。
与定理 \ref{thm:conclusion-phase-channel-crowding-lb} 矛盾，除非 \(\varepsilon\) 满足所示上界。证毕。
\end{proof}

\begin{definition}[\texorpdfstring{$U(1)$}{U(1)}-实现维数]\label{def:conclusion-ucdim}
对离散交换群 \(G\)，定义
\[
\mathrm{ucdim}(G):=\min\left\{k\in\NN:\ \exists\ \text{群单射 }G\hookrightarrow \TT^k\right\}.
\]
\end{definition}

\begin{proposition}[\texorpdfstring{$\mathrm{ucdim}(\ZZ^r)=1$}{ucdim(Z^r)=1}]\label{prop:conclusion-ucdim-zr}
对任意 \(r\ge 1\)，有 \(\mathrm{ucdim}(\ZZ^r)=1\)。
\leanverified{paper\_conclusion\_ucdim\_zr}
\end{proposition}

\begin{proof}
由定理 \ref{thm:conclusion-zn-to-circle-injective-dense} 存在群单射 \(\ZZ^r\hookrightarrow \TT\)，故 \(\mathrm{ucdim}(\ZZ^r)\le 1\)。
又 \(\TT^k\) 平凡群情形外不允许 \(k=0\)，故最小值为 \(1\)。证毕。
\end{proof}

\begin{proposition}[乘法自然数在单相位中的单射同态实现]\label{prop:conclusion-multiplicative-n-to-circle-mono-injection}
\leanverified{paper\_conclusion\_multiplicative\_n\_to\_circle\_mono\_injection}
令 \(\TT=\RR/\ZZ\)，并以自然对数定义映射
\[
\Phi:\NN^\times\to\TT,
\qquad
\Phi(n):=\log n\ \bmod 1.
\]
则 \(\Phi\) 为单射单子同态（乘法 \(\to\) 加法）：
\[
\Phi(ab)=\Phi(a)+\Phi(b),
\qquad
\ker\Phi=\{1\}.
\]
\end{proposition}

\begin{proof}
同态性由 \(\log(ab)=\log a+\log b\) 立即得到。
若 \(\Phi(n)=\Phi(m)\)，则 \(\log(n/m)\in\ZZ\)，即 \(n/m=e^k\)（某 \(k\in\ZZ\)）。
当 \(k\neq 0\) 时，由 Lindemann 定理 \(e^k\) 为超越数，从而 \(e^k\notin\QQ\)，与 \(n/m\in\QQ\) 矛盾；故 \(k=0\) 且 \(n=m\)。证毕。
\end{proof}

\begin{proposition}[正有理乘法群的显式单相位单射同态]\label{prop:conclusion-qpos-sqrtp-phase-injection}
\leanverified{paper\_conclusion\_qpos\_sqrtp\_phase\_injection}
以素数集 \(\mathcal P\) 为指标，按唯一分解有
\(
\QQ_{>0}^\times\cong \bigoplus_{p\in\mathcal P}\ZZ
\)。
定义
\[
\Psi:\QQ_{>0}^\times\to\TT,
\qquad
\Psi\!\left(\prod_{p}p^{e_p}\right):=\sum_{p}e_p\sqrt{p}\ \bmod 1,
\]
其中仅有限个 \(e_p\neq 0\)。
则 \(\Psi\) 为单射群同态。
\end{proposition}

\begin{proof}
同态性由指数加法立即推出。设 \(\sum_{i=1}^n c_i\sqrt{p_i}\in\ZZ\)（\(c_i\in\ZZ\)，\(p_i\) 互异素数）。
取域 \(K=\QQ(\sqrt{p_1},\dots,\sqrt{p_n})\)。对每个 \(j\) 存在自同构 \(\sigma_j\in\mathrm{Gal}(K/\QQ)\) 使
\(\sigma_j(\sqrt{p_j})=-\sqrt{p_j}\) 且对 \(i\neq j\) 有 \(\sigma_j(\sqrt{p_i})=\sqrt{p_i}\)。
对等式施以 \(\sigma_1\) 并与原式相减得 \(2c_1\sqrt{p_1}=0\)，故 \(c_1=0\)；依次得到 \(c_i=0\)。
因此 \(\sum_p e_p\sqrt p\in\ZZ\) 仅在全体 \(e_p=0\) 时成立，从而 \(\ker\Psi=\{1\}\)。证毕。
\end{proof}

\begin{theorem}[素数账本的单圈相位化：同态、单射与稠密像]\label{thm:conclusion-qpos-general-singlephase-injective-dense}
设 \(\mathbb P\) 为素数集，\(\Theta=\{\theta_p\}_{p\in\mathbb P}\subset\RR\) 满足
\[
\{1\}\cup\{\theta_p:\ p\in\mathbb P\}\ \text{在}\ \QQ\text{-线性意义下线性无关}.
\]
定义
\[
\Phi_\Theta:\QQ_{>0}^\times\to\TT,\qquad
\Phi_\Theta(q):=\exp\!\left(2\pi i\sum_{p\in\mathbb P}v_p(q)\theta_p\right),
\]
其中仅有限个 \(v_p(q)\neq 0\)。则：
\begin{enumerate}
  \item \(\Phi_\Theta\) 为群同态；
  \item \(\Phi_\Theta\) 为单射；
  \item \(\Phi_\Theta(\QQ_{>0}^\times)\) 在 \(\TT\) 中稠密。
\end{enumerate}
\leanverified{paper_conclusion_qpos_general_singlephase_injective_dense_package}
\end{theorem}

\begin{proof}
同态性由 \(v_p(q_1q_2)=v_p(q_1)+v_p(q_2)\) 直接得到。
若 \(\Phi_\Theta(q)=1\)，则存在 \(k\in\ZZ\) 使
\[
\sum_{p}v_p(q)\theta_p-k=0.
\]
这给出 \(\{1\}\cup\{\theta_p\}\) 的一个有理线性关系；由线性无关性得 \(k=0\) 且全部 \(v_p(q)=0\)，故 \(q=1\)，从而 \(\Phi_\Theta\) 单射。

由单射知像无限。又 \(\Phi_\Theta(\QQ_{>0}^\times)\) 为 \(\TT\) 的子群，按定理 \ref{thm:conclusion-circle-closed-subgroup-classification}，无限子群的闭包只能是 \(\TT\) 本身，故像稠密。证毕。
\end{proof}

\begin{remark}[单圈相位化的有限精度不可逆性]\label{rem:conclusion-qpos-singlephase-finite-precision-irrev}
定理 \ref{thm:conclusion-qpos-general-singlephase-injective-dense} 给出实现层的严格分界：\(\QQ_{>0}^\times\) 可单圈单射编码，但像在 \(\TT\) 中稠密。
因此任何固定分辨率的量化读出都会把无穷多个不同账本点压缩到同一相位桶，无法在全域上实现可逆解码。
\end{remark}

\begin{theorem}[多平方根线性型的范数下界]\label{thm:conclusion-multisqrt-norm-lower-bound}
\leanverified{paper\_conclusion\_multisqrt\_norm\_lower\_bound}
取两两不同素数 \(p_1,\dots,p_r\)（\(r\ge 1\)）。
对任意 \(a_1,\dots,a_r,k\in\ZZ\)，定义
\[
\alpha:=\sum_{j=1}^r a_j\sqrt{p_j}-k.
\]
若 \(\alpha\neq 0\)，则有
\[
|\alpha|
\ge
\left(|k|+\sum_{j=1}^r |a_j|\sqrt{p_j}\right)^{-(2^r-1)}.
\]
\end{theorem}

\begin{proof}
令 \(K=\QQ(\sqrt{p_1},\dots,\sqrt{p_r})\)。其 Galois 群同构于 \((\ZZ/2\ZZ)^r\)，每个自同构对应于独立改变各 \(\sqrt{p_j}\) 的符号。
由 \(\sqrt{p_j}\) 为代数整数知 \(\alpha\in\mathcal O_K\)，故范数
\[
N_{K/\QQ}(\alpha)
=
\prod_{\varepsilon\in\{\pm1\}^r}
\left(\sum_{j=1}^r \varepsilon_j a_j\sqrt{p_j}-k\right)
\]
为非零整数，从而 \(|N_{K/\QQ}(\alpha)|\ge 1\)。

设
\[
M:=|k|+\sum_{j=1}^r |a_j|\sqrt{p_j}.
\]
则每个共轭因子的绝对值都不超过 \(M\)，于是
\[
|N_{K/\QQ}(\alpha)|\le |\alpha|\,M^{2^r-1}.
\]
结合 \(|N_{K/\QQ}(\alpha)|\ge 1\) 即得结论。证毕。
\end{proof}

\begin{corollary}[单圈相位化的有限盒精度窗口]\label{cor:conclusion-prime-singlephase-precision-window-exp2r}
取两两不同素数 \(p_1,\dots,p_r\) 并定义
\[
\phi(e_1,\dots,e_r):=\exp\!\left(2\pi i\sum_{j=1}^r e_j\sqrt{p_j}\right)\in\TT.
\]
对 \(A\ge 1\)，记
\[
\mathcal E_{r,A}:=\{-A,\dots,A\}^r,\qquad
\delta_{r,A}:=\min_{\substack{e\neq e'\\ e,e'\in\mathcal E_{r,A}}}
\mathrm{dist}_{\TT}\bigl(\phi(e),\phi(e')\bigr),
\]
\[
B_{r,A}:=\left\lceil\log_2\frac{2}{\delta_{r,A}}\right\rceil.
\]
则有双侧界
\[
r\log_2(2A+1)\ \le\ B_{r,A}\ \le\
1+(2^r-1)\log_2\!\bigl(3Ar\sqrt{p_r}\bigr).
\]
特别地，有限盒上的证书化精度上界包含 \(2^r-1\) 指数项。
\end{corollary}

\begin{proof}
先证右侧上界。任取 \(e\neq e'\)，令 \(a_j=e_j-e'_j\) 并定义
\[
\Delta:=\left\|\sum_{j=1}^r a_j\sqrt{p_j}\right\|_{\RR/\ZZ}.
\]
取最近整数 \(k\in\ZZ\) 使
\[
\Delta=\left|\sum_{j=1}^r a_j\sqrt{p_j}-k\right|.
\]
由命题 \ref{prop:conclusion-qpos-sqrtp-phase-injection}，当 \(e\neq e'\) 时该差不可能为 \(0\)，故可应用定理 \ref{thm:conclusion-multisqrt-norm-lower-bound}。
由 \(|a_j|\le A\) 得
\[
|k|\le Ar\sqrt{p_r}+\frac12,\qquad
|k|+\sum_{j=1}^r |a_j|\sqrt{p_j}\le 3Ar\sqrt{p_r}.
\]
并由定理 \ref{thm:conclusion-multisqrt-norm-lower-bound} 得
\[
\Delta\ge (3Ar\sqrt{p_r})^{-(2^r-1)}.
\]
取 \(e\neq e'\) 的最小值得
\[
\delta_{r,A}\ge (3Ar\sqrt{p_r})^{-(2^r-1)},
\]
故
\[
B_{r,A}
\le
1+(2^r-1)\log_2(3Ar\sqrt{p_r}).
\]

再证左侧下界。点集 \(\phi(\mathcal E_{r,A})\) 的基数为 \((2A+1)^r\)。
把单位圆按弧长分成 \((2A+1)^r\) 个等分弧，抽屉原理给出两点弧距不超过 \((2A+1)^{-r}\)，即
\[
\delta_{r,A}\le (2A+1)^{-r}.
\]
因此
\[
B_{r,A}
=\left\lceil\log_2\frac{2}{\delta_{r,A}}\right\rceil
\ge r\log_2(2A+1).
\]
证毕。
\end{proof}

\begin{remark}[单圈相位化精度窗口的数值审计]\label{rem:conclusion-singlephase-prime-ledger-audit}
脚本 \texttt{scripts/exp\_conclusion\_singlephase\_prime\_ledger\_precision\_audit.py}
给出有限盒分离度、范数下界随机抽检与稠密性启发式审计，并生成
\texttt{artifacts/export/conclusion\_singlephase\_prime\_ledger\_precision\_audit.json}
与表 \ref{tab:conclusion_singlephase_prime_ledger_precision_audit}。
\begin{table}[H]
\centering
\scriptsize
\setlength{\tabcolsep}{5pt}
\caption{单相位素数账本在有限盒上的精度审计。$\delta_{r,A}$ 表示 $[-A,A]^r$ 上观测到的最小圆弧分离度，$\delta_{\mathrm{cert}}=(3Ar\sqrt{p_r})^{-(2^r-1)}$ 表示由多平方根范数界给出的证书化下界。}
\label{tab:conclusion_singlephase_prime_ledger_precision_audit}
\begin{tabular}{c c c c c c c c}
\toprule
$r$ & $A$ & $|\mathcal E_{r,A}|$ & $\delta_{r,A}$ & $\delta_{\mathrm{cert}}$ & $B_{\mathrm{lb}}$ & $B_{r,A}$ & $B_{\mathrm{ub}}$\\
\midrule
2 & 2 & 25 & 2.458e-02 & 1.114e-04 & 4 & 7 & 15\\
3 & 2 & 125 & 7.327e-04 & 5.844e-12 & 6 & 12 & 39\\
4 & 2 & 625 & 2.934e-04 & 9.091e-28 & 9 & 13 & 91\\
5 & 1 & 243 & 1.026e-03 & 2.509e-53 & 7 & 11 & 176\\
\midrule
\multicolumn{8}{c}{\textit{稠密性启发式（固定 $r=3$，统计 $[-L,L]^3$ 上最大圆弧间隙）：}}\\
\multicolumn{2}{c}{L=4} & \multicolumn{2}{c}{N=729} & \multicolumn{4}{c}{$\max\,\mathrm{gap}=5.840e-03$}\\
\multicolumn{2}{c}{L=8} & \multicolumn{2}{c}{N=4913} & \multicolumn{4}{c}{$\max\,\mathrm{gap}=7.327e-04$}\\
\multicolumn{2}{c}{L=12} & \multicolumn{2}{c}{N=15625} & \multicolumn{4}{c}{$\max\,\mathrm{gap}=2.487e-04$}\\
\bottomrule
\end{tabular}
\end{table}

\end{remark}

\begin{theorem}[圆群闭子群分类与无限子群稠密性]\label{thm:conclusion-circle-closed-subgroup-classification}
\leanverified{paper\_conclusion\_circle\_closed\_subgroup\_classification}
设 \(H\le \TT\) 为闭子群，则必有
\[
H\cong \ZZ/n\ZZ\ (n\in\NN)\qquad\text{或}\qquad H=\TT.
\]
因此任意无限子群 \(G\le\TT\) 在 \(\TT\) 中稠密。
\end{theorem}

\begin{proof}
记商映射 \(\pi:\RR\to\RR/\ZZ\simeq\TT\)，并置 \(\widetilde H:=\pi^{-1}(H)\)。
则 \(\widetilde H\) 为 \(\RR\) 的闭子群且 \(\ZZ\subseteq \widetilde H\)。
\(\RR\) 的闭子群仅有 \(a\ZZ\)（\(a\ge 0\)）与 \(\RR\) 两类。
若 \(\widetilde H=\RR\)，则 \(H=\TT\)。
若 \(\widetilde H=a\ZZ\)，由 \(\ZZ\subseteq a\ZZ\) 得 \(a=1/n\)（某 \(n\in\NN\)），从而
\[
H=\widetilde H/\ZZ\cong (1/n)\ZZ/\ZZ\cong \ZZ/n\ZZ.
\]
结论即得。证毕。
\end{proof}

\begin{theorem}[平方自由素数账本的指数级拥挤]\label{thm:conclusion-squarefree-crowding-kphase}
\leanverified{paper\_conclusion\_squarefree\_crowding\_kphase}
设 \(k\in\NN\) 且 \(\iota:\NN^\times\hookrightarrow \TT^k\) 为任意单射映射。
令 \(p_1,\dots,p_m\) 为前 \(m\) 个素数，并定义平方自由集合
\[
\mathcal S_m:=\left\{\prod_{j=1}^{m}p_j^{\varepsilon_j}:\ \varepsilon_j\in\{0,1\}\right\},
\qquad
|\mathcal S_m|=2^m.
\]
记最小分离度
\[
\delta_m:=\min_{\substack{x\neq y\\ x,y\in\mathcal S_m}} d_\infty\bigl(\iota(x),\iota(y)\bigr).
\]
则对所有 \(m\ge 1\) 有
\[
\delta_m\le \frac{1}{\max\{1,\lceil 2^{m/k}\rceil-1\}}
\ \lesssim\ 2^{-m/k}.
\]
\end{theorem}

\begin{proof}
将 \([0,1)^k\) 分割为 \(Q^k\) 个边长 \(1/Q\) 的半开立方体，其中 \(Q:=\max\{1,\lceil 2^{m/k}\rceil-1\}\)。
则 \(Q^k<2^m=|\mathcal S_m|\)（当 \(Q=1\) 时断言显然成立），故存在不同 \(x,y\in\mathcal S_m\) 使 \(\iota(x),\iota(y)\) 落于同一立方体，从而
\(
d_\infty(\iota(x),\iota(y))\le 1/Q
\)。
取最小即得所述不等式。证毕。
\end{proof}

\begin{corollary}[固定相位维数与固定分辨率下的可审计性破裂]\label{cor:conclusion-squarefree-audit-breakdown}
固定 \(k\in\NN\) 与目标分辨率 \(\varepsilon\in(0,1)\)。对任意单射映射 \(\iota:\NN^\times\hookrightarrow\TT^k\)，存在 \(m_0=m_0(k,\varepsilon)\) 使当 \(m\ge m_0\) 时，\(\iota\) 在 \(\mathcal S_m\) 上不可能保持 \(\varepsilon\)-分离。
更定量地，若要求对某个 \(m\) 在 \(\mathcal S_m\) 上满足
\(
d_\infty(\iota(x),\iota(y))\ge \varepsilon
\)
对所有 \(x\neq y\) 成立，则必有
\[
k\ \ge\ \frac{m}{\log_2(1/\varepsilon)}-O(1).
\]
\end{corollary}

% --- Distillation writeback ---
\begin{theorem}[有限相位桶的边界碰撞账本与外置严格化判别]\label{distill:grothendieck-compactified_boundary_degeneration_control-conclusion-phase-bucket-boundary-ledger-strictification}
令 \(G\) 为离散交换群，令 \(E\subset G\) 为非空有限审计族，并令
\[
\iota:E\longrightarrow \TT^k
\]
为任意相位读出。固定一个有限相位桶划分
\[
\mathcal B=\{B_\lambda\}_{\lambda\in\Lambda}
\]
并记
\[
\pi:E\to\Lambda,\qquad \pi(x)=\lambda\quad\Longleftrightarrow\quad \iota(x)\in B_\lambda.
\]
对每个非单点纤维
\[
C\in\mathcal C_\pi:=\{\pi^{-1}(\lambda):|\pi^{-1}(\lambda)|\ge2\}
\]
添入一个边界碰撞地层 \(\partial C\)，并定义其差分账本
\[
\Delta_C:=\{x-y:x,y\in C,\ x\ne y\}\subset G.
\]
则：
\begin{enumerate}
  \item 相位桶读出 \(\pi\) 在 \(E\) 上可逆，当且仅当 \(\mathcal C_\pi=\varnothing\)。若 \(\mathcal C_\pi\ne\varnothing\)，所有有限精度不可逆性都由边界地层族 \(\{\partial C\}_{C\in\mathcal C_\pi}\) 分类。
  \item 设 \(R\) 为交换群，\(\chi:G\to R\) 为授权外置群同态。则联合读出
  \[
  E\longrightarrow \Lambda\times R,\qquad x\longmapsto (\pi(x),\chi(x))
  \]
  在 \(E\) 上单射，当且仅当对每个 \(C\in\mathcal C_\pi\) 都有
  \[
  \ker\chi\cap\Delta_C=\varnothing.
  \]
  \item 若上述联合读出单射，则对每个 \(C\in\mathcal C_\pi\)，集合 \(\chi(C)\subset R\) 的基数为 \(|C|\)。特别地，任意有限外置寄存器都满足预算下界
  \[
  |R|\ge \max_{C\in\mathcal C_\pi}|C|.
  \]
\end{enumerate}
因此，有限相位精度下的失败不是单圈相位实现本身的矛盾，而是一个明确的边界碰撞账本；要消去它，外置寄存器必须在每个边界纤维上分离全部差分类。
\end{theorem}
\begin{proof}
第一项只是有限相位桶读出的精确分离条件。若每个纤维均为单点，则 \(\pi\) 显然在 \(E\) 上单射；若存在 \(|C|\ge2\)，取不同 \(x,y\in C\)，则 \(\pi(x)=\pi(y)\)，有限桶读出不可逆。把每个这样的非单点纤维记为一个边界地层 \(\partial C\)，恰好记录了全部被相位桶压成同一可见值的审计点。

证明第二项。若联合读出单射，取任意 \(C\in\mathcal C_\pi\) 与任意 \(d=x-y\in\Delta_C\)。若 \(d\in\ker\chi\)，则 \(\chi(x)=\chi(y)\)。又因 \(x,y\in C\)，有 \(\pi(x)=\pi(y)\)，从而联合读出不能区分 \(x\) 与 \(y\)，矛盾。因此 \(\ker\chi\cap\Delta_C=\varnothing\)。

反过来，设对每个边界纤维均有 \(\ker\chi\cap\Delta_C=\varnothing\)。若两个点 \(x,y\in E\) 的联合读出相等，则 \(\pi(x)=\pi(y)\)，故 \(x,y\) 属于同一 \(\pi\)-纤维。若该纤维为单点，则 \(x=y\)。若该纤维为某个 \(C\in\mathcal C_\pi\) 且 \(x\ne y\)，则
\[
x-y\in\Delta_C
\]
且 \(\chi(x-y)=0\)，这与 \(\ker\chi\cap\Delta_C=\varnothing\) 矛盾。因此必有 \(x=y\)，联合读出单射。

第三项中，若联合读出单射，则在固定纤维 \(C\) 上第一坐标 \(\pi\) 为常数，所以单射性完全由第二坐标 \(\chi|_C\) 承担。因此 \(\chi|_C:C\to R\) 为单射，\(|\chi(C)|=|C|\)，从而 \(|R|\ge |C|\)。对所有非单点纤维取最大值得到所示下界。

该判别只使用授权的群同态外置读出；它不把任意连续相位坐标误当作有限离散账本。故有限精度下的不可逆性被严格定位为边界纤维差分账本，外置化成本由这些差分类的分离需求给出。证毕。
\end{proof}

% --- End distillation writeback ---

\begin{proof}
由定理 \ref{thm:conclusion-squarefree-crowding-kphase} 有 \(\delta_m\le 1/(\max\{1,\lceil 2^{m/k}\rceil-1\})\)。
若 \(\varepsilon\)-分离成立，则 \(\delta_m\ge \varepsilon\)，从而 \(\varepsilon\le 1/(\lceil 2^{m/k}\rceil-1)\)，整理得所示下界。证毕。
\end{proof}

\begin{conjecture}[平方自由层上的最优条件数常数]\label{conj:conclusion-squarefree-optimal-conditioning}
对单射映射 \(\iota:\NN^\times\hookrightarrow \TT^k\) 定义
\[
C(\iota):=\liminf_{m\to\infty} 2^{m/k}\,\delta_m,
\]
其中 \(\delta_m\) 如定理 \ref{thm:conclusion-squarefree-crowding-kphase}。
预期对每个固定 \(k\) 存在常数 \(0<C_k^-\le C_k^+<\infty\) 使
\[
C_k^-\ \le\ \sup_{\iota} C(\iota)\ \le\ C_k^+,
\]
并且上确界可由算术上自然的相位同态实现达到。
\end{conjecture}

